

\usepackage[utf8]{inputenc} % Codifica caratteri
\usepackage[T1]{fontenc}
\usepackage[includehead, inner=3.0cm, outer=2.0cm, top=2.5cm, bottom=2.5cm]{geometry}
\usepackage{charter} % Text font
\usepackage[titles]{tocloft} % TOC compatta
\usepackage{amssymb} % Checkmark e altri simboli
\usepackage{graphicx} % Per inserire immagini
\usepackage{float} % Per controllare il posizionamento di immagini e tabelle
\usepackage[bottom]{footmisc} % Footer in fondo alla pagina
\usepackage[section]{placeins} % Force all floats before section end

\usepackage[english,italian]{babel}
\usepackage{amsmath, caption}
\usepackage[skip=8pt]{parskip}

\usepackage{listings}
\renewcommand{\lstlistingname}{Listato}
\usepackage{xcolor}
\usepackage{setspace}
\usepackage{csquotes}
\usepackage[
backend=biber,
style=apa,
sorting=ynt,
uniquelist=false]{biblatex}
\addbibresource{bibliography.bib} % Bibliografia su un file a parte
%\addbibresource{refs.bib} % File contentente la bibliografia


\let\Chapter\chapter\def\chapter{\addtocontents{lol}{\protect\addvspace{10pt}}\Chapter}

% Glossary
\usepackage[acronym,nonumberlist]{glossaries}
%\makeglossaries

% Liste
\usepackage{enumitem} % Per inserire liste
\usepackage{siunitx} % Per usare \setitemize
\setitemize{noitemsep,topsep=0pt,parsep=0pt,partopsep=0pt} % itemize più compatto
\setenumerate{noitemsep,topsep=0pt,parsep=0pt,partopsep=0pt} % enumerate più compatto
\setenumerate{noitemsep,topsep=0pt,parsep=0pt,partopsep=0pt} % enumerate più compatto
\setdescription{topsep=0pt} % enumerate più compatto

% Captions
\usepackage[font=small,textformat=period]{caption}
\captionsetup{justification=centering}

% Numerazione enumerate innestati
\renewcommand{\labelenumii}{\theenumii}
\renewcommand{\theenumii}{\theenumi.\arabic{enumii}.}

% Spaziature
\linespread{1.5} % Interlinea
\setlength{\parindent}{6pt} 
\setlength{\parskip}{1em} % Spazio tra paragrafi
\setlength{\headheight}{15pt} % Risolve warning con fancyhdr

% Formato capitoli
\usepackage{titlesec}
\titleformat{\chapter}[hang]{\normalfont\Huge\bfseries}{\thechapter.}{8pt}{\Huge\bfseries}
\titleformat*{\section}{\LARGE\bfseries}
\titleformat*{\subsection}{\Large\bfseries}
\titleformat*{\subsubsection}{\large\bfseries}
\titlespacing*{\chapter}{0pt}{-40pt}{15pt}
\titlespacing*{\section}{0pt}{18pt}{8pt}
\titlespacing*{\subsection}{0pt}{14pt}{6pt}
\titlespacing*{\subsubsection}{0pt}{10pt}{4pt}

% Intestazione pagine
\usepackage{fancyhdr}
\pagestyle{fancy}
\renewcommand{\chaptermark}[1]{\markboth{\thechapter. #1}{#1}}
\fancyhf{}
\fancyhead[L]{\small\slshape\nouppercase{\leftmark}}
\fancyfoot[C]{\thepage}
\renewcommand{\headrulewidth}{0.4pt}


\makeatletter
\renewcommand\paragraph{%
  \@startsection{paragraph}
    {4}
    {\z@}
    {3.25ex \@plus1ex \@minus.2ex}
    {-1em}
    {\normalfont\normalsize\bfseries\maybe@addperiod}%
}
\newcommand{\maybe@addperiod}[1]{%
  #1\@addpunct{.}%
}
\makeatother

\newcommand{\figura}[4]{%
  \begin{figure}[htbp]
    \centering
    \includegraphics[width=#1\textwidth]{#2}
    \caption{#3}\label{#4}
  \end{figure}%
}


% URL
\usepackage[hidelinks]{hyperref} % Per inserire URL
\usepackage{hyperxmp}

